% Authors: Carmen Beltran, Brendan Butler, Teddy Wachtler
%%%%%%%%%%%%%%%%%%%%%%%%%%%%%%%%%%%%%%%%%%%%%%%

\documentclass{article}
\usepackage[dvips]{graphicx}
\usepackage{a4wide}
\usepackage{amsmath}
\usepackage{euscript}
\usepackage{amssymb}
\usepackage{amsthm}
\usepackage{amsopn}

\theoremstyle{definition}
\newtheorem*{definition}{Definition}
\newtheorem{theorem}{Theorem}

\renewcommand{\qed}{\blacksquare}
\newcommand{\powset}{{\cal P}}
\newcommand{\img}{_{*}}
\newcommand{\pimg}{^{*}}

\newcommand{\vv}{\ensuremath{\vec{v}}}
\newcommand{\vu}{\ensuremath{\vec{u}}}
\newcommand{\vw}{\ensuremath{\vec{w}}}
\newcommand{\vx}{\ensuremath{\vec{x}}}
\newcommand{\vy}{\ensuremath{\vec{y}}}
\newcommand{\vb}{\ensuremath{\vec{b}}}
\newcommand{\vo}{\ensuremath{\vec{0}}}
\newcommand{\va}{\ensuremath{\vec{a}}}
\newcommand{\ve}{\ensuremath{\vec{e}}}

\newcommand{\R}{\mathbb{R}}
\newcommand{\Z}{\mathbb{Z}}
\newcommand{\C}{\mathbb{C}}
\newcommand{\N}{\mathbb{N}}
\newcommand{\Q}{\mathbb{Q}}

%%%%%%%%%%%%%%%%%%%%%%%%%%%%%%%%%%%%%%%%%%%%%%%

\begin{document}
	\begin{center}
		\Large{Math 251W: Foundations of Advanced Mathematics}
		\vspace{0.1cm}

		\normalsize Portfolio Assignment 4: \S 4.3 \& 4.4
		\vspace{0.1cm}

		\hfill Name: Carmen Beltran, Brendan Butler, Teddy Wachtler
		\vspace{0.1cm}

		\hrule
		\vspace{0.5cm}
	\end{center}

%%%%%%%%%%%%%%%%%%%%%%%%%%%%%%%%%%%%%%%%%%%%%%%

	\noindent Problem 4.3.10 [DISCUSS]
	\vspace{0.2cm}

	\underline{Proposition:}
	Let $f : A \rightarrow B$ be a function.

	\begin{itemize}
		\item[i]
		If $f$ has two distinct left inverses, it has no right inverse.

	\end{itemize}

	\vspace{0.2cm}
	\fbox{Proof} ( )
	\vspace{0.2cm}

	$\qed$.

	\vspace{0.5cm}

%%%%%%%%%%%%%%%%%%%%%%%%%%%%%%%%%%%%%%%%%%%%%%%

	\noindent Problem 4.3.11 [GROUP]
	\vspace{0.2cm}

	\underline{Proposition:}
	Let $f : A \rightarrow A_1 \times A_2 \times \dots \times A_k$ be a function, and let $U_i \subseteq A_i$ for each $i \in\{1, \ldots, k\}$, then $f \pimg (U_1 \times U_2 \times \dots \times U_k) = \bigcap_{i = 1}^k f_i \pimg (U_i)$ where $f_i$ is the $i^{th}$ coordinate function, $f_i = \pi_i \circ f$

	\vspace{0.2cm}
	\fbox{Proof} (Direct)
	\vspace{0.2cm}

	Let $k$ be an arbitrary natural number.
	Let $A, A_i$, $U_i$ be arbitrary sets so that $U_i \subseteq A_i$ for all $i \in \{1, \dots, k\}$.
	Let $f$ be an arbitrary function defined by $f : A \rightarrow A_1 \times A_2 \times \dots \times A_k$.
	Let $f_i$ be the coordinate function of $f$ defined by $f_i : A \rightarrow A_i$ and $f_i = \pi_i \circ f$ for all $i \in \{1, \dots, k\}$.
	We ask whether $f \pimg (U_1 \times U_2 \times \dots \times U_k) = \bigcap_{i = 1}^k f_i \pimg (U_i)$.

	Let $x$ be an arbitrary element of $f \pimg (U_1 \times U_2 \times \dots \times U_k)$.
	By the definition of the pre-image, we see that $f(x) \in U_1 \times U_2 \times \dots \times U_k$.
	By the definition of the Cartesian product, we see that $\pi_i(f(x)) \in U_i$ for all $i \in \{1, \dots, k\}$.
	By the definition of function composition, we see that $(\pi_i \circ f)(x) \in U_i$ for all $i \in \{1, \dots, k\}$.
	By the definition of $f_i$, we see that $f_i(x) \in U_i$ for all $i \in \{1, \dots, k\}$.
	By the definition of the pre-image, we see that $x \in f_i \pimg (U_i)$ for all $i \in \{1, \dots, k\}$.
	By the definition of indexed intersections, we see that $x \in \bigcap_{i = 1}^k f_i \pimg (U_i)$.
	Therefore, by the definition of subsets, since $x$ was an arbitrary element of $f \pimg (U_1 \times U_2 \times \dots \times U_k)$ and $x \in \bigcap_{i = 1}^k f_i \pimg (U_i)$, we see that $f \pimg (U_1 \times U_2 \times \dots \times U_k) \subseteq \bigcap_{i = 1}^k f_i \pimg (U_i)$.

	Let $y$ be an arbitrary element of $\bigcap_{i = 1}^k f_i \pimg (U_i)$.
	By the definition of indexed intersections, we see that $y \in f_i \pimg (U_i)$ for all $i \in \{1, \dots, k\}$.
	By the definition of the pre-image, we see that $f_i(y) \in U_i$ for all $i \in \{1, \dots, k\}$.
	By the definition of $f_i$, we see that $(\pi_i \circ f)(y) \in U_i$ for all $i \in \{1, \dots, k\}$.
	By the definition of function composition, we see that $\pi_i(f(y)) \in U_i$ for all $i \in \{1, \dots, k\}$.
	By the definition of the Cartesian product, we see that $f(y) \in U_1 \times U_2 \times \dots \times U_k$.
	By the definition of the pre-image, we see that $y \in f \pimg (U_1 \times U_2 \times \dots \times U_k)$.
	Therefore, by the definition of subsets, since $y$ was an arbitrary element of $\bigcap_{i = 1}^k f_i \pimg (U_i)$ and $y \in f \pimg (U_1 \times U_2 \times \dots \times U_k)$, we see that $\bigcap_{i = 1}^k f_i \pimg (U_i) \subseteq f \pimg (U_1 \times U_2 \times \dots \times U_k)$.

	By the definition of set equality, since $f \pimg (U_1 \times U_2 \times \dots \times U_k) \subseteq \bigcap_{i = 1}^k f_i \pimg (U_i)$ and $\bigcap_{i = 1}^k f_i \pimg (U_i) \subseteq f \pimg (U_1 \times U_2 \times \dots \times U_k)$, we see that $f \pimg (U_1 \times U_2 \times \dots \times U_k) = \bigcap_{i = 1}^k f_i \pimg (U_i)$.
	Therefore, for all natural numbers $k$, for all sets $A, A_i$, $U_i$ so that $U_i \subseteq A_i$ for all $i \in \{1, \dots, k\}$, and for all functions $f$ defined by $f : A \rightarrow A_1 \times A_2 \times \dots \times A_k$ and $f_i$ of $f$ defined by $f_i : A \rightarrow A_i$ and $f_i = \pi_i \circ f$ for all $i \in \{1, \dots, k\}$, we see that $f \pimg (U_1 \times U_2 \times \dots \times U_k) = \bigcap_{i = 1}^k f_i \pimg (U_i)$.
	$\qed$.

	\vspace{0.5cm}

%%%%%%%%%%%%%%%%%%%%%%%%%%%%%%%%%%%%%%%%%%%%%%%

	\noindent Problem 4.4.7 [DISCUSS]
	\vspace{0.2cm}

	\underline{Proposition:}
	The function $\phi : \powset (A) \rightarrow \powset (A)$, defined by $\phi (X) = A - X$ for all $X \in \powset (A)$, is bijective.

	\vspace{0.2cm}
	\fbox{Proof} ( )
	\vspace{0.2cm}

	$\qed$.

	\vspace{0.5cm}

%%%%%%%%%%%%%%%%%%%%%%%%%%%%%%%%%%%%%%%%%%%%%%%

	\noindent Problem 4.4.14 [INDIVIDUAL]
	\vspace{0.2cm}

	\begin{enumerate}
		\item[3]
		\underline{Proposition:}
		Given functions $f : A \rightarrow B$, and $g : B \rightarrow C$, if $g \circ f$ is bijective, then $f$ must be injective and $g$ must be surjective.

	\end{enumerate}

	\vspace{0.2cm}
	\fbox{Proof} (Direct)
	\vspace{0.2cm}

	Let $A, B$, and $C$ be arbitrary sets.
	Let $f$ and $g$ be arbitrary functions defined by $f : A \rightarrow B$ and $g : B \rightarrow C$.
	Suppose $g \circ f$ is bijective.
	By the definition of bijective functions, we see that $g \circ f$ is injective and surjective.

	Let $a_1$ and $a_2$ be arbitrary elements in $A$ so that $f(a_1) = f(a_2)$
	By substitution, we see that $g(f(a_1)) = g(f(a_2))$.
	By the definition of function composition, we see that $(g \circ f)(a_1) = (g \circ f)(a_2)$.
	By the definition of injective functions, we see that $a_1 = a_2$.
	By the definition of injective functions, since $a_1$ and $a_2$ were arbitrary elements of $A$, and $f(a_1) = f(a_2)$ implies $a_1 = a_2$, we see that $f$ is injective.

	Let $c$ be an arbitrary element in $C$.
	By the definition of surjective functions, since $g \circ f$ is surjective, we see that there exists $a \in A$ so that $(g \circ f)(a) = c$.
	By the definition of function composition, we see that $g(f(a)) = c$.
	By the definition of $f$, we see that there exists $b \in B$ so that $f(a) = b$.
	By substitution, we see that $g(b) = c$.
	Therefore, by the definition of surjective functions, since $c$ was an arbitrary element in $C$, and there exists an element $b \in B$ so that $g(b) = c$, we see that $g$ is surjective.

	Therefore, we see that for all sets $A, B$, and $C$, and for all sets $f : A \rightarrow B$, and $g : B \rightarrow C$, if $g \circ f$ is bijective, then $f$ must be injective and $g$ must be surjective.
	$\qed$.

	\vspace{0.5cm}

%%%%%%%%%%%%%%%%%%%%%%%%%%%%%%%%%%%%%%%%%%%%%%%
	\newpage

	\noindent Problem 4.4.17 [GROUP]
	\vspace{0.2cm}

	\underline{Proposition:}
	Let $f : A \rightarrow B$ be a map.
	$f$ is surjective iff $B - f \img (X) \subseteq f \img (A - X)$ for all $X \subseteq A$.

	\vspace{0.2cm}
	\fbox{Proof} (Direct)
	\vspace{0.2cm}

	Let $A$ and $B$ be arbitrary sets.
	Let $f$ be the function defined by $f : A \rightarrow B$.
	We ask whether $f$ is surjective if and only if $B - f \img (X) \subseteq f \img (A - X)$ for all sets $X \subseteq A$.

	Consider the case where $f$ is surjective.
	Let $X$ be an arbitrary set so that $X \subseteq A$.
	Let $y$ be an arbitrary element of $B - f \img (X)$.
	By the definition of differences of sets, we see that $y \in B$ and $y \not \in f \img (X)$.
	By the definition of surjective functions, since $y \in B$ and $f : A \rightarrow B$, there exists a particular element $x \in A$ so that $f(x) = y$.
	By the definition of images, since $y \not \in f \img (X)$, we see that for all $z \in X, f(z) \neq y$.
	Since $f(x) = y$, we see that $x \not \in X$.
	By the definition of differences of sets, since $x \in A$ and $x \not \in X$, we see that $x \in A - X$.
	By the definition of images, since $f(x) = y$ and $x \in A - X$, we see that $y \in f \img (A - X)$.
	By the definition of subsets, since $y$ was an arbitrary element of $B - f \img (X)$ and $y \in f \img (A - X)$, we see that $B - f \img (X) \subseteq f \img (A - X)$.
	Therefore, we see that if $f$ is surjective, then $B - f \img (X) \subseteq f \img (A - X)$ for all sets $X \subseteq A$.

	Consider the case where $B - f \img (X) \subseteq f \img (A - X)$ for all sets $X \subseteq A$.
	Take the case where $X = \emptyset$.
	By substitution, we see that $B - f \img (\emptyset) \subseteq f \img (A - \emptyset)$.
	By the definition of images, we see that $f \img (\emptyset) = \emptyset$.
	By the definition of differences of sets, we see that $A - \emptyset = A$.
	By substitution, we see that $B \subseteq f \img (A)$.
	Let $b$ be an arbitrary element of $B$.
	By the definition of subsets, we see that $b \in f \img (A)$.
	By the definition of images, there exists a particular element $a \in A$ so that $f(a) = b$.
	By the definition of surjective functions, since $b$ was an arbitrary element of $B$, we see that $f$ is surjective.
	Therefore, we see that if $B - f \img (X) \subseteq f \img (A - X)$ for all sets $X \subseteq A$, then $f$ is surjective.

	Since if $f$ is surjective, then $B - f \img (X) \subseteq f \img (A - X)$ for all $X \subseteq A$, and if $B - f \img (X) \subseteq f \img (A - X)$ for all $X \subseteq A$, then $f$ is surjective, we see that $f$ is surjective if and only if $B - f \img (X) \subseteq f \img (A - X)$ for all $X \subseteq A$.
	Therefore, for all sets $A$ and $B$, if $f$ is a function defined by $f : A \rightarrow B$, then $f$ is surjective if and only if $B - f \img (X) \subseteq f \img (A - X)$ for all $X \subseteq A$.
	$\qed$.

	\vspace{0.5cm}

%%%%%%%%%%%%%%%%%%%%%%%%%%%%%%%%%%%%%%%%%%%%%%%

	\noindent Challenge Problem
	\vspace{0.2cm}

	\underline{Proposition:}
	Let $f: A \rightarrow B$ be a function.
	We can define a new function $F: \mathcal{P}(A) \rightarrow \mathcal{P}(B)$ by the image map:\[ F(S) = f_{*}(S) = \{ f(x)  \mid x \in S \} \]

	\begin{enumerate}
		\item
		Prove that if $f$ is surjective, then $F$ is surjective.

		\vspace{0.2cm}
		\fbox{Proof} ( )
		\vspace{0.2cm}

		$\qed$.

		\vspace{0.2cm}
		\item
		Is it true that if $f$ is injective, then $F$ is injective?
		Provide a Proof or counterexample.

		\vspace{0.2cm}
		\fbox{Proof} ( )
		\vspace{0.2cm}

		$\qed$.

		\vspace{0.2cm}
	\end{enumerate}

%%%%%%%%%%%%%%%%%%%%%%%%%%%%%%%%%%%%%%%%%%%%%%%

\end{document}